\chapter{犯罪构成}

% Source PDF physical page 26; see audit/source_coverage.json.
\begin{EditionOpeningPanel}
新中国成立后，我国刑法学开始借鉴苏联版四要件体系，取得重大进步。七十多年后的今天，苏联早已解体，而我国的法治事业取得巨大成就。能够取得如此巨大成就，主要经验是，将马克思法治理论中国化、时代化。同理，应将苏联版四要件体系中国化、时代化，形成一种新时代改良版的四要件体系。这种犯罪构成体系吸收了阶层式体系的有益理念。
\end{EditionOpeningPanel}

\section{定罪标准：改良版四要件犯罪构成体系}

\subsection*{一、改良版四要件体系}

1.犯罪客体：法益

2.犯罪主体：（1）自然人；（2）单位

3.客观要件：（1）行为；（2）结果；（3）因果关系

4.主观要件：（1）故意；（2）过失

5.排除犯罪事由I-排除违法事由：

（1）正当防卫；（2）紧急避险；（3）被害人承诺

6.排除犯罪事由II-排除责任事由：

（1）责任年龄；（2）责任能力；（3）违法性认识可能性；（4）期待可能性

关于改良版四要件体系的解读：（亲爱的朋友，本节内容请认真读两遍）

1.犯罪客体

这是指犯罪行为所侵害的、法律所保护的客体。这个“客体”是指法律保护的利益，简称法益。例如，刑法设立盗窃罪，是为了保护国民的财产权。这个财产权就是犯罪行为所侵害的、刑法所保护的客体，被称为犯罪客体、保护客体或保护法益。

[注意］犯罪客体不是指犯罪对象（行为对象）。例如，甲盗窃到乙的手机。犯罪对象是乙的手机，犯罪客体是乙的财产权。犯罪对象是犯罪客体的载体，犯罪客体是犯罪对象蕴含的利益。

% Source PDF physical page 27; see audit/source_coverage.json.
2.犯罪主体\footnote{苏联版四要件体系将“犯罪主体”放在“客观要件”后面。然而，从事情的发生逻辑顺序看，应先有个犯罪主体，再实施犯罪行为，也即先有人，再有人的行为。}

犯罪主体也称行为主体。第一种犯罪主体是具有侵害法益能力的自然人。例如，10岁的甲盗窃他人手机，甲具有侵害法益的能力，就是合格的犯罪主体。第二种犯罪主体就是单位。

3.客观要件与主观要件

犯罪主体具备客观要件和主观要件，意味着制造了法益侵害事实，简称违法事实，其行为便具有法益侵害性和应被禁止性，合称违法性。例如，甲故意用刀杀死仇人乙。甲具备客观要件（杀人行为、死亡结果及因果关系），具备主观要件（杀人故意）。这意味着甲故意制造了侵害生命法益的事实，简称违法事实。甲的行为具有法益侵害性和应被禁止性，合称违法性。

4.排除犯罪事由I：排除违法事由\footnote{狭义的四要件体系不包括排除犯罪事由，而广义的四要件体系包括排除犯罪事由。}

行为人的行为具有违法性后，进一步审查发现，行为人存在某些事由，可以排除行为具有的违法性（在此指法益侵害性），简称排除违法事由，也称为正当化事由。例如，甲用刀杀死乙。事后查明，乙先用刀杀甲，甲出于正当防卫而杀死乙。甲的正当防卫能够排除“甲杀乙”所具有的违法性（法益侵害性）。正当防卫就是排除违法事由。排除违法事由也被称为正当化事由，意指行为被正当化了。如果事后查明，乙没有杀甲，乙在买水果，甲出于感情纠纷，用刀杀死乙，那么甲不存在正当防卫，最终确定甲制造了违法事实，甲的行为具有违法性。\footnote{认定一个行为具有违法性，包含两层含义，一是在实然层面，认为该行为具有法益侵害性；二是在应然层面，认为该行为具有应被禁止性。有些行为具有法益侵害性，但不具有刑法的禁止性。例如，排队时插队，具有法益侵害性，但不具有刑法的禁止性，不构成犯罪，这是因为刑法是最严厉的制裁措施，要保持谦抑性。排除违法事由中的违法性是指法益侵害性。后文“法律认识错误”中有个“违法性认识可能性”，其中的“违法性”是指刑法的禁止性。}

5.排除犯罪事由II：排除责任事由

\CJKunderwave{刑法的定罪机制是：判断坏事一判断坏人。先判断，有没有人干坏事，也即是否制造}\CJKunderwave{违法事实；后判断，这个人是不是个坏人，也即应否被刑法谴责}。

传统观念是，干坏事的人一定是坏人。实则不然。例如，甲盗窃了他人的手机，后查明甲只有10岁（盗窃罪的责任年龄是16岁），或者是精神病患者（缺乏责任能力）。甲缺乏法规范意识，亦即不知道自已行为的法律性质及法律后果。此时，刑法无法谴责甲，甲最终不用负刑事责任。这一点也符合我国古代法律思想。孔子曰：不教而诛谓之虐，不戒视成谓之暴。这句话用现代刑法理论诠释，意思就是，不教育一个人，使他有法规范意识，那么他做错事，就杀死他，这属于虐杀；不告诫一个人，使他能认识到自已行为的性质，而只检视成果，看后果是不是危害后果，这属于暴政。

因此，10岁的甲偷手机，属于干了坏事，但不是坏人。精神正常的成年人偷手机，属于干了坏事，是个坏人。精神病患者打人，属于干了坏事，但不是坏人。（脑子瓦特了的人，不算坏人，只能算病人）

责任年龄、责任能力等是谴责一个人的条件，也即评价一个人是不是“坏人”的条件。缺乏这些条件，就缺乏可谴责性，由此不能追究刑事责任。刑法将“可谴责性”简称

% Source PDF physical page 28; see audit/source_coverage.json.
为“有责性”或“责任”，将排除可谴责性的事由称为排除责任事由，也称为宽恕事由。这里的“责任”是指谴责（是个动词），不是指刑事责任（是个名词）。\footnote{刑法中的“责任”具有四种含义。第一种含义是“义务”。例如，“你有责任去救小孩”，意指“你有义务去救小孩”。第二种含义是“结果归责”，意指将结果归责给谁，谁应为结果负责。例如，“造成这个车祸，是你的责任”，意指“车祸是你造成的，要算到你的头上，由你负责”。能够将结果归责给一个人的前提条件是，这个人的行为与结果具有因果关系。第三种含义就是此处所指的“谴责”，是一种否定性价值评价。第四种含义是“刑事责任”，意指一种不利的法律后果。“承担刑事责任”是在具备可谴责性后得出的结论，是在谴责的基础上采取的惩罚措施。因此，在讨论“责任”时，应注意所使用的语境和含义，避免指称不清。}

[总结］犯罪构成体系的判断任务

第一，四个要件（客体、主体、客观要件、主观要件）的判断任务是，行为人有没有干坏事？

第二，排除违法事由的判断任务是，行为人干的事，到底是不是坏事？有没有可能是正当的事？

第三，排除责任事由的判断任务是，行为人是不是坏人？能不能谴责他？有没有可能宽恕他？\footnote{提示：以上犯罪构成体系是完整的“零部件清单”，不是每个犯罪的成立都需要具备这些全部的“零部件”，具体犯罪是按需搭配。例如，就故意杀人罪的成立而言，并不需要具备实害结果和因果关系。具备实害结果和因果关系是故意杀人罪的既遂条件。}

\subsection*{二、“犯罪”概念的阶段化理解}

对“犯罪”的含义，应当根据发展阶段进行理解，而不能理解成一个终局意义的“点”。“犯罪”是个阶段性的概念，而非“点”概念。

1.\CJKunderwave{“制造违法事实”意义上的“犯罪”概念：干了坏事；此时，对事不对人}

行为人具备客观要件和主观要件，且不具有排除违法事由，就表明行为人制造了法益侵害事实（违法事实），其行为具有法益侵害性（违法性）。这种“违法行为”就是一种“犯罪”行为，亦即“制造违法事实”意义上的犯罪行为。这是一种阶段化的犯罪概念。

2.\CJKunderwave{最终负刑事责任的“犯罪”概念：干了坏事，并且是坏人；此时，对事又对人}

如果行为人不具有排除责任事由，则这种阶段化的“犯罪”概念就转变成应负刑事责任的“犯罪”概念。如果行为人具有排除责任事由，则最终不负刑事责任；但是，不能由此认为，“行为人没有制造违法事实”。行为人仍存在“制造违法事实”意义上的犯罪行为。

例如，10岁的甲盗窃到乙的手机。甲的行为是“制造违法事实”意义上的“犯罪行为”。甲未达责任年龄，无法谴责甲，甲不负刑事责任，最终无罪。但是，不能因此否定甲的行为是“制造违法事实”意义上的“犯罪行为”。甲盗窃的手机属于“犯罪所得”。丙若为甲保管该手机，则构成掩饰、隐瞒犯罪所得罪。\footnote{参见2025年8月26日最高人民法院、最高人民检察院《关于办理掩饰、隐瞒犯罪所得、犯罪所得收益刑事案件适用法律若干问题的解释》第十条规定。}

\subsection*{三、法考真题答案采取改良版四要件体系}

真题（2012年第9题，15岁抢夺案），甲（15周岁）求乙（16周岁）为其抢夺作接应，乙同意。某夜，甲抢夺被害人的手提包（内有1万元现金），将包扔给乙，然后吸引被害人跑开。乙害怕坐牢，将包扔在草丛中，独自离去。关于本案，下列哪一选项是错误的？

% Source PDF physical page 29; see audit/source_coverage.json.
A.甲不满16周岁，不构成抢夺罪

B.甲与乙构成抢夺罪的共犯

C.乙不构成抢夺罪的间接正犯

D.乙成立抢夺罪的中止犯

官方答案：D项。根据中国化、改良版四要件体系，甲乙构成“制造违法事实”意义上的共同犯罪，甲是实行犯，乙是帮助犯。在排除责任（谴责）阶段，由于甲未达抢夺罪要求的16周岁，最终不负刑事责任；由于乙已满16周岁，最终负抢夺罪（帮助犯）的刑事责任。如果问实行犯在哪里，就是甲，“制造违法事实”意义上的实行犯。

苏联版四要件体系认为，犯罪主体必须是具有责任年龄、责任能力的人。依此，由于抢夺罪的责任年龄是16周岁，所以甲不是合格的犯罪主体，直接无罪。如此，对乙只能单独处理。单独处理，就对乙无法定罪，出现处罚漏洞。（1）对乙不能定帮助犯（也即共犯），是因为帮助犯不能单独存在。（2）对乙不能定间接正犯，是因为间接正犯要求对被利用人具有支配力，将其作为犯罪工具加以支配。本案是甲请求乙做自已的帮助犯。乙没有指使支配甲。本题答案也认为乙不构成间接正犯。可见，苏联版四要件体系会导致处罚漏洞。这是其不足之处。

按照苏联版四要件体系，A项正确，B项错误。然而，官方答案认为B项说法也正确。这是因为，苏联四要件体系将“犯罪”概念理解成一个终局意义的“点”概念，而按照改良版四要件体系，B项是“制造违法事实”意义上的判断，A项是进入排除责任（谴责）阶段的判断。这是阶段化理解“犯罪”概念。

\CJKunderwave{这表明，官方答案的立场是中国化、改良版的四要件体系。“九大本”法考教材刑法}\CJKunderwave{卷2026版第59页也支持上述改良版四要件体系的分析。}

\subsection*{四、“先客观、后主观”的审查顺序}

例如，甲在沙漠里看到前方站着仇人乙，向其开枪，实际是稻草人，四下无人。（暂不考虑枪支类犯罪，下同）

主观主义认为（少数说）：首先，甲具有杀人的犯罪故意，符合了主观要件；其次，判断客观要件，甲有开枪杀人行为，符合了客观要件，那么甲便构成犯罪；最后，由于没有杀死人，所以成立故意杀人罪（未遂）。

客观主义认为（多数说）：首先判断客观要件，核心是判断是否存在危害行为。要存在危害行为，就必须对法益产生实害或危险。如果一个行为对法益不会造成任何危险，就属于日常生活行为。本案中，甲有开枪行为，但开枪行为不一定就是危害行为。开枪行为要成为危害行为，必须对他人的生命产生威胁。甲现在开枪打的是稻草人，四下又无人，对他人的生命不会产生任何危险，所以不是危害行为。连危害行为都不是，对甲应直接得出不构成故意杀人罪的结论。

有人可能会想，那甲主观上有杀人故意啊！但是\CJKunderwave{注意}：犯罪是行为，而不是思想！只有杀人故意，而没有杀人的危害行为，不构成犯罪。刑法惩罚的是危害行为，而非思想。\footnote{提示：主观主义不是指定罪时只看主观，不看客观；客观主义也不是指定罪时只看客观，不看主观。主观主义以主观优先，客观主义以客观优先。}

% Source PDF physical page 30; see audit/source_coverage.json.
\section{定罪方法：三段论推理}

[问题］某孤儿院领导们开会决定，将孤儿小芳卖掉，小芳被卖到一个青楼。刑法规定，拐卖儿童罪只能由自然人构成，单位不构成本罪。对本案该如何处理？（孤儿被卖案，2011年第2题）

关于三段论推理，举例说明，大前提：凡是人都会死；小前提：柏拉图是人；结论：柏拉图会死。三段论推理应用到法律中，法律规定是大前提，案件事实是小前提。根据小前提是否符合大前提，得出有罪或无罪的结论。刑法学的主要任务有两个：一是如何解释大前提（法律规定），也即刑法条文规定的犯罪构成要件；二是如何认定小前提（案件事实）。\footnote{正确理解美国霍姆斯大法官的名言“法律的生命不在于逻辑，而在于经验”。霍姆斯并不反对逻辑的重要性，只是为了批判当时盛行的一种观点“法律发展的唯一动力在于逻辑”。该名言的准确表述应是“法律的生命不仅在于逻辑，还在于经验”。英美法系侧重经验主义，采取判例法。大陆法系侧重理性主义，采取成文法。逻辑与经验都很重要，不能偏废，二者之间有个度的平衡问题。我国是大陆法系背景，对于我国的法律初学者而言，在学习法律初期，在尚不具备逻辑推理能力、尚未建立逻辑体系时，应当理性主义先行，打造自己的逻辑推理能力，在此过程中逐渐积累经验。因此需要正确理解霍姆斯的这句话的时代背景。当你对逻辑推理已经掌握得非常娴熟，甚至唯逻辑论时，霍姆斯的这句话对你是有意义的。}

\InsertDiagram{figures/manual_p030_1.tex}

\subsection*{一、大前提}

\subsubsection*{（一）大前提的范围}

判断有罪无罪的大前提只能是刑法规定的犯罪构成要件。\CJKunderwave{有罪无罪的唯一标准是犯罪}\CJKunderwave{构成要件}。刑法分则为每个罪名都规定了犯罪构成要件；符合了，就有罪，不符合，就无罪；除此之外，再无别的标准。

当然，完整的大前提是犯罪构成体系，行为符合四个要件后，还要看有无排除犯罪事\CJKunderwave{由（排除违法事由和排除责任事由）}。

[提示】道德不能成为有罪无罪的判断标准。刑事法庭不是道德法庭，刑事审判不是道德审判。有观点认为，符合道德的行为一定不构成犯罪。这种看法有误。例如，以前发生一个案件，儿子是乡里恶霸，欺男霸女，无恶不作，没人管得了，老父亲大义灭亲，将儿子毒死。若依据这种看法，老父亲能够大义灭亲，道德非常高尚，因此不构成犯罪。这种结论显然错误。老父亲肯定构成犯罪，只是量刑上可以从宽处罚。又如，电影《周处除三害》中，男主角为民除害，杀死逍遥法外的大恶人，该行为符合道德、道义，但仍构成故意杀人罪。道德的归道德，法律的归法律。

\subsubsection*{（二）大前提的种类}

大前提（刑法分则规定的具体罪状）由一个个构成要件要素组成。这些要素可以分为以下种类：

% Source PDF physical page 31; see audit/source_coverage.json.
1.记述的构成要件要素和规范的构成要件要素

\CJKunderwave{记述的构成要件要素}，是指只需根据客观上的事实判断即可确定的要素。“记述”是指描述，描述事实。例如，“伪造或变造”“护照、签证”。

\CJKunderwave{规范的构成要件要素}，是指需要法官根据主观上的价值判断才能确定的要素。例如，“猥亵”“侮辱”“淫秽物品”。这类要素的特点是，受人的主观价值观的影响较大。例如，对于一幅画，观念保守的人可能会认为这是一件淫秽物品，而观念开放的人可能会认为这是一幅艺术作品，正所谓仁者见仁，“淫”者见“淫”。\CJKunderwave{注意}：这里的“规范”的含义，不是指法律规范、刑事规范等规范的意思，而是指哲学上价值评价的意思。（2008年第51题）

[提示】\CJKunderwave{“事实判断”与“价值评价”}是刑法学经常用到的重要概念。二者区别在于，对事实判断的结论可以作真假判断，例如，“前方站着一个男人”，这是事实判断，对该结论可以作真假判断。对价值评价的结论无法作真假判断。例如，“前方这个男人是个坏男人”，这是价值评价，对该结论无法作真假判断，因为每个人的价值观不同，结论也就不同。

2.成文的构成要件要素和不成文的构成要件要素

成文的构成要件要素，是指刑法条文明文规定的要素。例如，刑法第236条明文规定强奸罪的对象是“妇女”，第263条明文规定抢劫罪的手段是“暴力、胁迫或者其他方法”。

不成文的构成要件要素，是指刑法条文表面上没有规定，但实质上必须具备的要素。

例如，刑法第264条规定的盗窃罪、第266条规定的诈骗罪都没有规定“以非法占有为目的”，但实际上这些财产犯罪必须“以非法占有为目的”。“以非法占有为目的”就是这些罪名的不成文的构成要件要素。

3.积极的构成要件要素和消极的构成要件要素

积极的构成要件要素，是指积极的、正面的表明犯罪成立的要素。例如，徇私枉法罪中的“司法工作人员”便是积极的构成要件要素，只有司法工作人员才能构成徇私枉法罪。（2014年第4题）

消极的构成要件要素，是指消极的、反面的否定犯罪成立的要素。例如，刑法第389条第3款行贿罪中规定：“因被勒索给予国家工作人员以财物，没有获得不正当利益的，不是行贿。”这是从反面否定了行贿罪的成立，属于消极的构成要件要素。（目前公认的消极的构成要件要素就这一个例子）

4.客观的构成要件要素和主观的构成要件要素

客观的构成要件要素，是指表明行为外在的、客观方面的要素。例如，行为、结果、时间、地点等。

主观的构成要件要素，是指表明行为人内心的、主观方面的要素。例如，故意、过失、目的、动机等。

\subsection*{二、小前提}

认定小前提主要是刑事诉讼法及其证据法的任务。刑法所要处理的是确实无法查明的情形。对于确实无法查明的情形，基于保障人权，应适用存疑时有利于被告的原则。

\subsubsection*{（一）适用领域}

该原则不是解释大前提（法律规定）的原则，而是解决悬疑事实的原则。在事实不

% Source PDF physical page 32; see audit/source_coverage.json.
清、存在疑问时，应根据存疑时有利于被告的原则认定事实。在解释法律含义时，不应该也不可能一味按照存疑时有利于被告的原则来解释，否则法益保护无从谈起。

例如，甲杀死仇人乙后，发现乙的钱包，便拿走钱包。案件事实很清楚，但甲的取财行为符合盗窃罪的法律规定还是符合侵占罪的法律规定，存在疑问。对此，需要对理论深人研究，而不能根据存疑时有利于被告的原则一味定侵占罪。如果该原则能够适用于解释法律规定，则给了法官询私枉法的尚方宝剑，法官完全可以将上述被告人的行为解释成无罪，这样最有利于被告人。

\subsubsection*{（二）存疑的情形}

案件事实存疑，是指案件事实存在两种或三种可能性。对此只能择一认定，选择一个有利于被告人的事实加以认定。

1.要么有罪、要么无罪，认定为无罪。

在调查案件事实时，如果既存在有罪事实的可能，也存在无罪事实的可能，根据存疑时有利于被告的原则，应当认定为无罪事实。例如，甲的轿车丢失，警察在乙的住所找到。证据能够证明，乙要么替小偷窝藏了该车，要么不知情而购买到赃车。对此，应认定为无罪。

2.要么重罪、要么轻罪，认定为轻罪。

如果既存在重罪事实的可能，也存在轻罪事实的可能，不存在无罪事实的可能，根据存疑时有利于被告的原则，应当认定为轻罪事实。例如，甲的轿车丢失，警察在乙的住所找到。证据能够证明，乙要么亲自盗窃了该车，要么替他人窝藏了该车。对此，对乙应认定为窝藏赃物。

3.要么重罪、要么轻罪、要么无罪，认定为无罪。

例如，甲的轿车丢失，警察在乙的住所找到。证据能够证明，乙要么亲自盗窃了该车，要么替小偷窝藏了该车，要么不知情而购买到赃车。对此，应认定为无罪。

\subsection*{三、正确推导}

\subsubsection*{（一）符合性的判断与案件事实的提炼}

小前提（案件事实）符合大前提（构成要件），则行为构成犯罪。对案件事实可以从不同角度进行提炼归纳。

例如，对于本节开头“孤儿被卖案”，有种推导过程是：案件事实是单位在实施拐卖儿童，而刑法规定，拐卖儿童罪只能由自然人构成，单位不构成本罪，所以，本案只能作无罪处理。

正确的推导是，对“单位实施拐卖儿童”这个案件事实进行提炼，其中必然存在“某些自然人（直接责任人员）在实施拐卖儿童”的事实。这些自然人制造的案件事实符合拐卖儿童罪的构成要件，因此这些自然人构成该罪。

\subsubsection*{（二）循环往复推导与想象竞合}

1.循环往复推导与全面审查。对同一个案件事实，需要从不同侧面、不同角度提炼案件事实，然后根据不同的大前提（构成要件）进行推导，有可能构成多个犯罪。

例如，杨某故意杀害正在执行公务的执法人员。一方面，从杀人这个角度提炼出“杨

% Source PDF physical page 33; see audit/source_coverage.json.
某故意杀人”的案件事实。对此，杨某构成故意杀人罪。另一方面，从妨害公务这个角度提炼出“杨某妨害公务”的案件事实。对此，杨某构成妨害公务罪。

2.想象竞合。这是循环往复使用三段论所造成的现象。一个行为（小前提，案件事实）同时触犯两个罪名。例如，杨某的一个行为同时触犯了故意杀人罪和妨害公务罪，两罪想象竞合，择一重罪论处，最终以故意杀人罪论处。

（1）特征：一个行为，触犯了两个或数个罪名。（一石二鸟）

（2）处理办法：择一重罪论处。这是因为，只有一个行为，如果数罪并罚，就意味着对一个行为进行了重复处罚，这是不公平的，侵犯了被告人的人权。

（3）明示功能。虽然对一个行为不能重复处罚，但是必须对一个行为全面审查，看这一个行为都触犯了哪些犯罪，制造了多少法益侵害事实。例如，必须在判决中指出来，上述杨某的一个行为同时触犯故意杀人罪和妨害公务罪，制造了两个法益侵害事实，一个是剥夺了他人的生命，另一个是妨害了正常的公务活动。这便是\CJKunderwave{想象竞合的明示功能}，也即明白显示这一个行为制造了多少个法益侵害事实。如果在判决书中，只指出杨某构成故意杀人罪，则会遗漏评价该行为制造的另一个法益侵害事实（妨害公务）。指出来，但不并罚。这种做法是在保护法益与保障人权之间的一个平衡做法。

[提示】想象竞合犯是对一个行为重复评价，但是没有重复处罚，而是择一重罪论处。\CJKunderwave{所谓“禁止重复评价原则”，实际是指禁止重复处罚}。如果重复处罚，则侵犯了犯罪人的人权。

\OriginalPanel{典型真题}

关于构成要件要素，下列哪一选项是错误的？（2014年·卷二·4题）\footnote{[答案]D。}

A.传播淫秽物品罪中的“淫秽物品”是规范的构成要件要素、客观的构成要件要素

B.签订、履行合同失职被骗罪中的“签订、履行”是记述的构成要件要素、积极的构成要件要素

C.“被害人基于认识错误处分财产”是诈骗罪中的客观的构成要件要素、不成文的构成要件要素

D.“国家工作人员”是受贿罪的主体要素、规范的构成要件要素、主观的构成要件要素

